%% Generated by doc/gen-options.py from `lassip --help`. Do not edit.
% 29 options in 6 sections.

\subsection{General}

\begin{description}[style=nextline,leftmargin=0pt,font=\normalfont\ttfamily]
\item[--threads <int>]
The number of threads to spawn during computations.
\\\textit{Default:} \texttt{1}
\item[--seed <int>]
Seed for the shuffle used when clustering haplotypes that carry missing data. Runs with the same seed, inputs and flags are reproducible regardless of \texttt{--threads}, operating system or compiler. Set 0 to seed from the system clock instead, which reproduces the non-reproducible behaviour of lassip <= 1.2.2.
\\\textit{Default:} \texttt{1}
\end{description}

\subsection{Input and output}

\begin{description}[style=nextline,leftmargin=0pt,font=\normalfont\ttfamily]
\item[--out <string>]
The basename for all output files.
\\\textit{Default:} \texttt{outfile}
\item[--map <string>]
A map file formatted <chr\#> <locusID> <genetic pos> <physical pos>. Sites in VCF not in map file will be interpolated.
\\\textit{Default:} \texttt{(empty)}
\item[--vcf <string1> ... <stringN>]
One or more VCF files containing haplotype data. Variants should be coded 0/1 and every file must contain all of the samples named by \texttt{--pop}. Several contigs may be analysed in one run, given either as one file per contig or as files holding several; the run then writes a single spectra file per population covering all of them. A contig may not appear twice, and a file's records must be grouped by contig.
\\\textit{Default:} \texttt{(empty)}
\item[--pop <string>]
A file containing <ind ID> <pop ID>.
\\\textit{Default:} \texttt{(empty)}
\item[--spectra <string1> ... <stringN>]
A list of spectra files for finalization. Contigs are delimited by the chr column rather than by the file, so a file may hold any number of them and the results are the same either way. A file's rows must be grouped by contig and ascending within one, and no contig may appear twice across the files given here.
\\\textit{Default:} \texttt{(empty)}
\item[--unphased <bool>]
Set this flag to indicate data are unphased.
\\\textit{Default:} \texttt{false}
\end{description}

\subsection{Windows}

\begin{description}[style=nextline,leftmargin=0pt,font=\normalfont\ttfamily]
\item[--winsize <int>]
The window size within which to calculate statistics.
\\\textit{Default:} \texttt{0}
\item[--winstep <int>]
The sliding window step size.
\\\textit{Default:} \texttt{0}
\end{description}

\subsection{Statistics}

\begin{description}[style=nextline,leftmargin=0pt,font=\normalfont\ttfamily]
\item[--calc-spec <bool>]
Set this flag to compute K-truncated haplotype frequency spectra.
\\\textit{Default:} \texttt{false}
\item[--avg-spec <bool>]
Set this flag to compute and output the average K-truncated haplotype frequency spectrum from a set of .spectra files.
\\\textit{Default:} \texttt{false}
\item[--null-spec <string>]
A file containing a null K-truncated haplotype spectrum for use computing LASSI or saltiLASSI.
\\\textit{Default:} \texttt{(empty)}
\item[--lassi <bool>]
Set this flag to use the LASSI method.
\\\textit{Default:} \texttt{false}
\item[--lassi-choice <int>]
Set this flag to change the way LASSI distributes mass across sweeping haplotype classes. Takes an integer in \{1..5\}.
\\\textit{Default:} \texttt{4}
\item[--hapstats <bool>]
Set this flag to calculate haplotype statistics.
\\\textit{Default:} \texttt{false}
\item[--salti <bool>]
Set this flag to use the saltiLASSI method.
\\\textit{Default:} \texttt{false}
\item[--k <int>]
Top K haplotypes for LASSI computations.
\\\textit{Default:} \texttt{10}
\end{description}

\subsection{Filtering}

\begin{description}[style=nextline,leftmargin=0pt,font=\normalfont\ttfamily]
\item[--filter-level <int>]
Filter monomorphic sites and sites with missing data: 0-no filtering, 1-compute freq for all samples, 2-compute freq per pop.
\\\textit{Default:} \texttt{2}
\item[--max-lmiss <double>]
Filter loci with > this proportion of missing data.
\\\textit{Default:} \texttt{0.10}
\item[--max-hmiss <double>]
Drop haplotypes with > this proportion of missing data when computing the HFS.
\\\textit{Default:} \texttt{0.20}
\item[--hap-cluster <string>]
How haplotypes carrying missing genotypes are grouped into classes. Missing sites act as wildcards, so such a haplotype can be compatible with several distinct haplotypes and something has to choose. One of: best-comp (default) each haplotype joins the most frequent class it is compatible with, taking the best-observed haplotypes first. Deterministic; \texttt{--seed} has no effect on it. garud-shuffle the pre-1.3 behaviour: shuffle the haplotypes, then let each in turn absorb every compatible one not yet claimed. The result depends on the shuffle, so runs differ unless \texttt{--seed} is fixed. soft-em EXPERIMENTAL. Divide an ambiguous haplotype's count across the classes it could belong to in proportion to their frequencies, iterated to convergence. Class sizes become fractional. Without missing genotypes and at \texttt{--match-tol} 0 all three give the same classes.
\\\textit{Default:} \texttt{best-comp}
\item[--match-tol <int>]
Group haplotypes with missing data into the same class as a haplotype with no missing data if they have <= this many pairwise differences.
\\\textit{Default:} \texttt{0}
\item[--keep-monomorphic <bool>]
Set this flag to retain any monomorphic sites in the data.
\\\textit{Default:} \texttt{false}
\end{description}

\subsection{saltiLASSI}

\begin{description}[style=nextline,leftmargin=0pt,font=\normalfont\ttfamily]
\item[--dist-type <string>]
Distance measure for saltiLASSI: bp, cm, nw.
\\\textit{Default:} \texttt{bp}
\item[--max-gap <double>]
Widest gap in basepairs between two genetic map positions that \texttt{--dist-type} cm will interpolate across. A window falling in a wider gap cannot be placed and lassip stops rather than guess its genetic position. Real maps have gaps at centromeres and other low-recombination regions, so raise this if your map is sparse where your data are dense. Set 0 to interpolate across any gap.
\\\textit{Default:} \texttt{3000000.00}
\item[--max-extend-bp <double>]
Maximum distance in basepairs from core window to consider for saltiLASSI.
\\\textit{Default:} \texttt{100000.00}
\item[--max-extend-cm <double>]
Maximum distance in centimorgans from core window to consider for saltiLASSI.
\\\textit{Default:} \texttt{0.05}
\item[--max-extend-nw <double>]
Maximum distance in number of windows from core window to consider for saltiLASSI.
\\\textit{Default:} \texttt{5.00}
\end{description}

